% Generated by scripts/assemble.py; corrections belong in source/corrections.json.
\examhead{21}{中国人民大学 2026 年数学分析试卷}{25}
% p026-b10; source PDF page 26
\blockanchor{p026-b10}
\question{1}(15 分) 讨论 $f(x) = \sin \sqrt{x}$ 和 $g(x) = \sin (x^2)$ 在 $[0, +\infty)$ 上是否一致连续.\par
% p026-b11; source PDF page 26
\blockanchor{p026-b11}
\question{2}(15 分) 设 $a_{n} > 0, n = 1, 2, \cdots$, 且 $\lim_{n \to \infty} \frac{a_{n+1}}{a_n} < +\infty$, 证明: $\lim_{n \to \infty} \sqrt[n]{a_n} \leq \lim_{n \to \infty} \frac{a_{n+1}}{a_n}$.\par
% p026-b12; source PDF page 26
\blockanchor{p026-b12}
\question{3}(15 分) 证明: 对任意的 $x \in [0,1]$, 有 $1 + x^2 \leq 2^x$.\par
% p026-b13; source PDF page 26
\par\Needspace{6\baselineskip}
\blockanchor{p026-b13}
\question{4}(15 分) 设 $f(x)$ 为 $\mathbb{R}$ 上周期为 $\pi$ 的函数, 且在 $\left[-\frac{\pi}{2}, \frac{\pi}{2}\right)$ 上, 有\par
\nopagebreak[4]
% p026-b14; source PDF page 26
\blockanchor{p026-b14}
\displayformula{f (x) = \left\{ \begin{array}{l l} \frac {\pi}{2} + x, & x \in \left[ - \frac {\pi}{2}, 0\right); \\ \frac {\pi}{2} - x, & x \in \left[ 0, \frac {\pi}{2}\right). \end{array} \right.}
% p026-b15; source PDF page 26
\blockanchor{p026-b15}
请将 $f(x)$ 在 $\left[-\frac{\pi}{2}, \frac{\pi}{2}\right)$ 上展开为傅里叶级数, 并由此求 $\sum_{n=1}^{\infty} \frac{1}{(2n-1)^2}$.\par
% p026-b16; source PDF page 26
\blockanchor{p026-b16}
\question{5}(15 分) 设对于任意的正整数 $n$, $f_n(x)$ 均为定义在 $[a, b]$ 上的黎曼可积函数, 且 $\{f_n(x)\}$ 在 $[a, b]$ 上一致收敛于函数 $f(x)$. 证明: $f(x)$ 是 $[a, b]$ 上的黎曼可积函数.\par
% p026-b17; source PDF page 26
\blockanchor{p026-b17}
\question{6}(15 分) 计算积分 $\int_{0}^{\frac{\pi}{2}} \frac{1}{\sin \theta} \ln \frac{1 + t \sin \theta}{1 - t \sin \theta} \mathrm{d}\theta, |t| < 1.$\par
% p026-b18; source PDF page 26
\par\Needspace{6\baselineskip}
\blockanchor{p026-b18}
\question{7}(15 分) 设\par
\nopagebreak[4]
% p026-b19; source PDF page 26
\blockanchor{p026-b19}
\displayformula{f (x, y) = \left\{ \begin{array}{l l} \frac {x ^ {2} y + x ^ {4}}{x ^ {2} + y ^ {2}}, & x ^ {2} + y ^ {2} \neq 0; \\ 0, & x ^ {2} + y ^ {2} = 0. \end{array} \right.}
% p026-b20; source PDF page 26
\blockanchor{p026-b20}
讨论 $f(x,y)$ 在 $(0,0)$ 的连续性和可微性. 又设 $x(t) = t, y(t) = t^2$, 讨论复合函数 $f(x(t),y(t))$ 在 $t = 0$ 的导数是否满足链式法则.\par
% p027-b00; source PDF page 27
\par\Needspace{6\baselineskip}
\blockanchor{p027-b00}
\question{8}(15 分) 设 $f(x, y)$ 为定义在 $D = \{(x, y) \mid x^2 + y^2 < 1\}$ 上的二元函数, 在 $D$ 内对 $x$ 和 $y$ 具有直到二阶的连续偏导数, 且 $(x_0, y_0) \in D$ 是 $f(x, y)$ 的极小值点. 证明:\par
\nopagebreak[4]
% p027-b01; source PDF page 27
\blockanchor{p027-b01}
\displayformula{f _ {x x} (x _ {0}, y _ {0}) + 2 f _ {x y} (x _ {0}, y _ {0}) + f _ {y y} (x _ {0}, y _ {0}) \geq 0.}
% p027-b02; source PDF page 27
\par\Needspace{6\baselineskip}
\blockanchor{p027-b02}
\question{9}(15 分) 设反常积分\par
\nopagebreak[4]
% p027-b03; source PDF page 27
\blockanchor{p027-b03}
\displayformula{I = \iiint_ {D} \frac {x ^ {2}}{(x ^ {2} + y ^ {2} + z ^ {2}) ^ {2}} \mathrm{d} x \mathrm{d} y \mathrm{d} z.}
% p027-b04; source PDF page 27
\blockanchor{p027-b04}
其中 $D = \{(x,y,z)\mid x^2 +y^2 +z^2\leq 4x\}$ .证明反常积分 $I$ 收敛，并求出 $I$ 的值.\par
% p027-b05; source PDF page 27
\blockanchor{p027-b05}
\question{10}(15 分) 设 $D$ 是平面区域, $u(x, y)$ 为定义在 $D$ 上的二元函数, 且在 $D$ 内对 $x$ 和 $y$ 具有直到二阶的连续偏导数. 证明: $u(x, y)$ 在 $D$ 上满足 $\frac{\partial^2 u}{\partial x^2} + \frac{\partial^2 u}{\partial y^2} = 0$ 的充要条件为对于任意 $D$ 内的圆周 $L$, 且 $L$ 所围圆 $O$ 含于 $D$, 有 $\oint_L \frac{\partial u}{\partial n} \mathrm{d}s = 0$, 其中 $\pmb{n}$ 取圆周 $L$ 的外法向.\par
